% ============================================================
\chapter*{Preface}
\addcontentsline{toc}{chapter}{Preface}
% ============================================================

\section*{Course Objectives}

This \textbf{Stochastic Processes} course is aimed at Master's-level students in
mathematics, theoretical physics, quantitative finance, and engineering sciences.
It provides a rigorous, measure-theoretic presentation of the main classes of
random processes and their applications.

\medskip

The primary objective is to guide the student from probability foundations through
to the construction of the It\^o stochastic integral and the solution of stochastic
differential equations (SDEs), via in-depth study of Markov chains, the Poisson
process, martingales, and Brownian motion.

\medskip

Upon completing this course, students will be able to:
\begin{itemize}
  \item Master the measure-theoretic formalism of stochastic processes.
  \item Analyse Markov chains in discrete and continuous time: classify states,
    compute stationary distributions, establish ergodicity.
  \item Work with martingales: apply fundamental inequalities, convergence
    theorems, and the optional stopping theorem.
  \item Understand the fine properties of Brownian motion.
  \item Construct the It\^o stochastic integral and apply It\^o's formula.
  \item Solve stochastic differential equations and apply them in finance,
    physics, and biology.
\end{itemize}

\section*{Prerequisites}

The essential prerequisites are:
\begin{itemize}
  \item \textbf{Measure theory and integration}: $\sigma$-algebras, measures,
    Lebesgue integral, convergence theorems (dominated, monotone), $L^p$ spaces.
  \item \textbf{Undergraduate probability}: random variables, expectation,
    variance, law of large numbers, central limit theorem.
  \item \textbf{Elementary functional analysis}: Banach and Hilbert spaces,
    modes of convergence, completeness.
  \item \textbf{Linear algebra}: matrices, eigenvalues, spectral decomposition.
  \item \textbf{Ordinary differential equations}: existence and uniqueness
    theorems, solution methods.
\end{itemize}

\section*{Course Organisation}

The course is organised into eleven chapters following a logical progression:

\begin{enumerate}
  \item \textbf{Probability Review}: probability spaces, conditional expectation,
    modes of convergence. This chapter provides the measure-theoretic foundations
    needed for the rest of the course.

  \item \textbf{Stochastic Processes --- Definitions}: formal definition of a
    stochastic process, filtrations, stopping times, adapted and predictable
    processes. These notions form the basic language.

  \item \textbf{Discrete-Time Markov Chains}: Markov property, state
    classification, ergodic theorems, stationary measures.

  \item \textbf{Continuous-Time Markov Chains}: infinitesimal generator,
    Kolmogorov equations, birth-and-death processes.

  \item \textbf{Poisson Process}: axiomatic construction, fundamental properties,
    compound and non-homogeneous Poisson processes.

  \item \textbf{Martingales --- Discrete Time}: definition, fundamental
    inequalities (Doob), convergence theorems, optional stopping theorem.

  \item \textbf{Martingales --- Continuous Time}: path regularity, Doob--Meyer
    decomposition, local martingales.

  \item \textbf{Brownian Motion}: construction (L\'evy--Ciesielski), path
    properties, nowhere differentiability, law of the iterated logarithm,
    reflection principle.

  \item \textbf{It\^o Stochastic Integral}: construction via simple processes,
    isometric extension to $L^2$ processes, martingale properties.

  \item \textbf{It\^o's Formula and SDEs}: It\^o change-of-variables formula,
    multidimensional It\^o formula, existence and uniqueness of strong solutions.

  \item \textbf{Applications}: finance (Black--Scholes model, hedging), physics
    (Langevin equation, diffusion), biology (stochastic population models).
\end{enumerate}

\section*{Conventions and Notation}

Throughout this course, we adopt the following conventions:

\begin{itemize}
  \item $(\Omega, \mathcal{F}, \PP)$ denotes a fixed probability space.
  \item $\{\mathcal{F}_t\}_{t \geq 0}$ or $\{\mathcal{F}_n\}_{n \geq 0}$ denotes
    a filtration satisfying the usual conditions (completeness, right-continuity).
  \item $\E[X]$ denotes the expectation of the random variable $X$.
  \item $\E[X \mid \mathcal{G}]$ denotes the conditional expectation of $X$ given
    the $\sigma$-algebra $\mathcal{G}$.
  \item $\PP(A)$ denotes the probability of the event $A$.
  \item $\Var(X)$ and $\Cov(X, Y)$ denote variance and covariance.
  \item $\R$, $\N$, $\Z$, $\C$ denote the sets of reals, naturals, integers,
    and complex numbers. We write $\R_+ = [0, +\infty)$.
  \item $\norm{\cdot}$ and $\abs{\cdot}$ denote a norm and the absolute value.
  \item $\mathbf{1}_A$ denotes the indicator function of the set $A$.
  \item a.s.\ means ``almost surely'', r.v.\ means ``random variable'',
    SDE means ``stochastic differential equation''.
  \item $\xrightarrow{\text{a.s.}}$, $\xrightarrow{L^p}$, $\xrightarrow{\PP}$,
    $\xrightarrow{\mathcal{L}}$ denote almost sure, $L^p$, in probability, and
    distributional convergence.
\end{itemize}

\section*{Pedagogical Approach}

Each chapter contains:
\begin{itemize}
  \item Precise \textbf{definitions} within the measure-theoretic framework.
  \item \textbf{Theorems} with complete or sketched proofs.
  \item Detailed \textbf{examples} illustrating abstract concepts.
  \item \textbf{Exercises} of varying difficulty, some with hints.
  \item \emph{Intuition} boxes to develop probabilistic intuition.
  \item \emph{Warning} boxes to flag common pitfalls.
  \item \emph{Key formulas} boxes summarising essential results.
\end{itemize}

\section*{Bibliographic References}

The following texts have inspired this course:

\begin{itemize}
  \item \textsc{Br\'emaud}, P. \emph{Markov Chains: Gibbs Fields, Monte Carlo
    Simulation and Queues}. Springer, 1999.
  \item \textsc{Durrett}, R. \emph{Probability: Theory and Examples}. Cambridge
    University Press, 5th ed., 2019.
  \item \textsc{Karatzas}, I. and \textsc{Shreve}, S.E. \emph{Brownian Motion and
    Stochastic Calculus}. Springer, 2nd ed., 1991.
  \item \textsc{Le Gall}, J.-F. \emph{Brownian Motion, Martingales, and Stochastic
    Calculus}. Springer, 2016.
  \item \textsc{Norris}, J.R. \emph{Markov Chains}. Cambridge University Press, 1997.
  \item \textsc{{\O}ksendal}, B. \emph{Stochastic Differential Equations}. Springer,
    6th ed., 2003.
  \item \textsc{Revuz}, D. and \textsc{Yor}, M. \emph{Continuous Martingales and
    Brownian Motion}. Springer, 3rd ed., 1999.
  \item \textsc{Shreve}, S.E. \emph{Stochastic Calculus for Finance~I \& II}.
    Springer, 2004.
  \item \textsc{Williams}, D. \emph{Probability with Martingales}. Cambridge
    University Press, 1991.
\end{itemize}

\section*{Acknowledgements}

The author thanks colleagues and students whose remarks and suggestions have
helped improve this course over the years. Discussions with practitioners in
quantitative finance and physicists have also enriched the presentation of
applications.

\vfill

\begin{flushright}
\textit{The Author}\\
March 2026
\end{flushright}
